\begin{lemma}
\label{lemma-homothety-spectrum}
Let $R \to S$ be a ring map.
Let $N$ be a $S$-module.
Assume
\begin{enumerate}
\item $R$ is a local ring with maximal ideal $\mathfrak m$,
\item $\overline{S} = S/\mathfrak m S$ is Noetherian, and
\item $\overline{N} = N/\mathfrak m_R N$ is a finite $\overline{S}$-module.
\end{enumerate}
Let $\Sigma \subset S$ be the multiplicative subset of elements which are not
a zerodivisor on $\overline{N}$. Then $\Sigma^{-1}S$ is a semi-local ring
whose spectrum consists of primes $\mathfrak q \subset S$ contained in an
element of $\text{Ass}_S(\overline{N})$. Moreover, any maximal
ideal of $\Sigma^{-1}S$ corresponds to an associated prime of
$\overline{N}$ over $\overline{S}$.
\end{lemma}

\begin{proof}
Note that
$\text{Ass}_S(\overline{N}) = \text{Ass}_{\overline{S}}(\overline{N})$, see
Algebra, Lemma \ref{algebra-lemma-ass-quotient-ring}.
This is a finite set by
Algebra, Lemma \ref{algebra-lemma-finite-ass}.
Say $\{\mathfrak q_1, \ldots, \mathfrak q_r\} = \text{Ass}_S(\overline{N})$.
We have $\Sigma = S \setminus (\bigcup \mathfrak q_i)$ by
Algebra, Lemma \ref{algebra-lemma-ass-zero-divisors}.
By the description of $\Spec(\Sigma^{-1}S)$ in
Algebra, Lemma \ref{algebra-lemma-spec-localization}
and by
Algebra, Lemma \ref{algebra-lemma-silly}
we see that the primes of $\Sigma^{-1}S$ correspond to the primes of
$S$ contained in one of the $\mathfrak q_i$.
Hence the maximal ideals of $\Sigma^{-1}S$ correspond one-to-one with the
maximal (w.r.t.\ inclusion) elements of the set
$\{\mathfrak q_1, \ldots, \mathfrak q_r\}$. This proves the lemma.
\end{proof}
